\Article{14}
Ist also die Funktion \(y\) auf die Form einer Summe von Gliedern wie
\[
Ax^\alpha x'{}^\beta x''{}^\gamma\mathord{.}\,\mathord{.}\,\mathord{.}\,\mathord{.}
\]
gebracht, so wird der Wert des Integrals \(\int y\psi y.\mathrm{d}y\) bei Erstreckung über alle Werte von \(y\), also der Mittelwert von \(y\), gleich der Summe der Glieder
\[
A\times\int x^\alpha\varphi x.\mathrm{d}x\times\int x'{}^\beta\varphi x'.\mathrm{d}x'\times\int x''{}^\gamma\varphi x''.\mathrm{d}x''\mathord{.}\,\mathord{.}\,\mathord{.}\,\mathord{.}
\]
sein, wobei die Integrationen von \(x=-\infty\) bis \(x=+\infty\), von \(x'=-\infty\) bis \(x'=+\infty\) etc. zu erstrecken sind; oder, was auf dasselbe hinausläuft, gleich der Summe der Glieder, die entstehen, wenn für die einzelnen Potenzen \(x^\alpha,x'{}^\beta,x''{}^\gamma\) etc. ihre Mittelwerte eingesetzt werden. Die Wahrheit dieses sehr wichtigen Satzes hätte sich auch aus anderen Überlegungen leicht herleiten lassen.

\Article{15}
Wenden wir das im vorigen Artikel Dargelegte auf den besonderen Fall an, in dem
\[
y=\frac{xx+x'x'+x''x''+\text{etc.}}{\sigma}
\]
ist, wobei \(\sigma\) die Anzahl der Glieder im Zähler bezeichnet. Der Mittelwert von \(y\) ergibt sich hier sofort \(=mm\), wenn das Zeichen \(m\) in derselben Bedeutung wie oben genommen wird. Der tatsächliche Wert von \(y\) kann in einem bestimmten Fall freilich größer oder kleiner als der Mittelwert ausfallen, ebenso wie der tatsächliche Wert des einzelnen Gliedes \(xx\): Die Wahrscheinlichkeit aber, dass der zufällige Wert von \(y\) nicht merklich vom Mittelwert \(mm\) abweicht, wird sich mit wachsender Anzahl \(\sigma\) immer mehr der Gewissheit nähern. Um dies deutlicher zu machen, wollen wir, da es nicht in unserer Macht steht, die Wahrscheinlichkeit selbst genau zu bestimmen, den zu befürchtenden mittleren Fehler untersuchen, wenn wir \(y=mm\) annehmen. Nach den Grundsätzen von Art. 6 wird dieser Fehler offensichtlich die Quadratwurzel aus dem Mittelwert der Funktion
\[
\left(\frac{xx+x'x'+x''x''+\text{etc.}}{\sigma}-mm\right)^2
\]
sein. Um diesen zu ermitteln, genügt es zu bemerken, dass der Mittelwert eines Gliedes wie \(\frac{x^4}{\sigma\sigma}\) \(=\frac{n^4}{\sigma\sigma}\) ist (wenn das Zeichen \(n\) in der Bedeutung von Art. 11 verwendet wird), der Mittelwert eines Gliedes \SourceLeaf{020}wie \(\frac{2xxx'x'}{\sigma\sigma}\) dagegen \(=\frac{2m^4}{\sigma\sigma}\) wird, woraus sich der Mittelwert jener Funktion sehr leicht ergibt
\[
=\frac{n^4-m^4}{\sigma}
\]

Hieraus lernen wir, dass sich, wenn eine genügend große Menge voneinander unabhängiger zufälliger Fehler \(x,x',x''\) etc. vorliegt, daraus mit großer Gewissheit ein Näherungswert von \(m\) durch die Formel
\[
m=\BareRoot{\frac{(xx+x'x'+x''x''+\text{etc.})}{\sigma}}
\]
gewinnen lässt und dass der bei dieser Bestimmung zu befürchtende mittlere Fehler in Bezug auf das Quadrat \(mm\)
\[
=\BareRoot{\frac{n^4-m^4}{\sigma}}
\]
ist. Da die letztere Formel die Größe \(n\) enthält, wird es, wenn nur irgendeine Vorstellung vom Grad der Genauigkeit dieser Bestimmung gewonnen werden soll, genügen, eine Hypothese über die Funktion \(\varphi\) anzunehmen. Zum Beispiel wird bei der dritten Hypothese von Art. 9, 11 jener Fehler \(=mm\BareRoot{\frac2\sigma}\). Falls dies weniger zusagt, kann ein Näherungswert von \(n^4\) aus den Fehlern selbst mithilfe der Formel
\[
\frac{x^4+x'{}^4+x''{}^4+\text{etc.}}{\sigma}
\]
gewonnen werden. Allgemein können wir aber behaupten, dass die doppelte Genauigkeit bei dieser Bestimmung die vierfache Menge an Fehlern erfordert oder dass das Gewicht der Bestimmung der Anzahl \(\sigma\) proportional ist.

Ganz ebenso lässt sich, wenn die Beobachtungsfehler einen konstanten Teil enthalten, dessen Näherungswert umso sicherer aus dem arithmetischen Mittel vieler Fehler erschließen, je größer deren Anzahl ist. Dabei wird der bei dieser Bestimmung zu befürchtende mittlere Fehler durch
\[
\BareRoot{\frac{mm-kk}{\sigma}}
\]
ausgedrückt, wenn \(k\) den konstanten Teil selbst und \(m\) den mittleren Fehler der noch nicht vom konstanten Teil befreiten Beobachtungen bezeichnet, oder einfach durch \(\frac{m}{\BareRoot{\sigma}}\), wenn \(m\) den mittleren Fehler der vom konstanten Teil befreiten Beobachtungen bezeichnet (s. Art. 8).

\clearpage
